\explainedfigure{fig_headline_success_clear_weather_200_episodes}{0.97}{Success rate of the driving models on the same 200 test episodes}{Each bar is the share of 200 identical test episodes (scenario seeds 5000 to 5199, clear weather, sparse traffic) in which a model reached the destination without a collision and without leaving the road. The black line is the 95\% Wilson confidence interval, the range of success rates that is statistically compatible with the observed count. Grey bars are rule-based teachers that see the exact simulator state; orange bars are the two baseline models supplied by the teammates; blue bars are models trained in this project.}{The Proposed Bird's-Eye-View Driving Model succeeds in 92.5\% of the episodes. The two baseline models succeed in 85.0\% (trained without traffic) and 74.5\% (trained with traffic). The proposed model is about as good as the Collision-Aware Teacher (95.0\%) whose behaviour it learned; the intervals overlap, so the two cannot be told apart. The Image-Only Driving Model, which reads camera pixels instead of the bird's-eye-view grid, reaches 0.130 (interval 0.090 to 0.184), far below every model that uses the grid.}{Training a model by letting it drive and asking a teacher what it should have done (dataset aggregation) produces a model that matches its teacher and clearly beats few-demonstration behaviour cloning. This figure does not say \emph{which part} of the proposed design (the training method or the map overlay) causes the gain; that needs the additional evaluations listed in the paper-directions section.}{fig:headline}
